\cotefile{121}
\scanpage{49}
\unsure{Contextes}
\[
\left\{\begin{array}{l}
\mathrm{L.P.M.}\\
\mathrm{S.A.P.M.}\\
\mathrm{Alg}_{\mathbb C},\;\mathrm{An}_{\mathbb C}
\end{array}\right.
\]

Soient $X$ donné, avec une famille loc. finie de sous-espaces fermés $T_\alpha\subset X$. Alors il existe une partition \inserted{\uuline{strictement}} admissible (\deleted{\unsure{strictement}} liée aux conditions de connexité) de $X$, pour laquelle les $T_\alpha$ sont réunions de \unsure{feuilles}. Parmi celles-ci, il y en a une plus grossière que \unsure{ttes|$\forall$ les} autres. Sa formation commute à la restriction à \unsure{tt|un} ouvert.
\margin{\uline{NB} c'est une \unsure{propriété} \uline{locale} sur $X$}

La canonicité de cette dernière montre qu'il suffit de prouver l'existence loc$^{\mathrm t}$, et la propriété \unsure{d'induction}. On peut supposer donc $X$ de dim finie, \deleted{\unsure{et} $X$ \illegible{3}} \inserted{\unsure{Nbre} des $T_\alpha$ fini. On procède par réc. sur dim $X$.} Trivial si $\dim X\unsure{=}0$ (unique partition strict$^{\mathrm t}$ adm. \unsure{obt.})

\deleted{Considérons les}

Supposons $\dim X>0$.

On va définir \inserted{(ce qui sera) la réunion} des \uline{feuilles ouvertes} de la partition cherchée, \inserted{\unsure{en sera un premier ouvert}} \deleted{Considérons \unsure{tous} les \unsure{intersect.} des $T_\alpha$, \unsure{c. à. d.} ($r\geq1$ facteurs) $T_{\alpha_1}\cap T_{\alpha_2}\cap\cdots\cap T_{\alpha_r}$ qui sont} l'ouvert $X^{\unsure{1}}$, comme
\[
X^{\unsure{1}}=\bigcup_\alpha \mathring{T}_\alpha
\]

Soit

\margin{50}
